\section{The full spin-core exclusion}
\label{sec:spin-quarter-obstruction}

The complete local spectra sharpen HD-SPIN-01 and remove both
its self-duality and its nonsplit-polarization hypotheses.

\begin{lemma}[HD-SPIN-02a: the local half-eigenspace bound]
\label{HD-SPIN-02a}
Let $U=2.\mathfrak A_6$ and let $W$ be any complex representation
on which the central involution acts as $-1$.  For every
$h\in U\setminus Z(U)$ and every eigenvalue $\lambda$,
\[
                   \dim\ker(h-\lambda)\le\tfrac12\dim W.
                                                               \tag{SPINQ.1}
\]
\end{lemma}

\begin{proof}
It suffices to check each irreducible constituent, since the
multiplicities of a fixed eigenvalue add.  There are six
center-negative rows in the complete published table $2.A6$.
The following table gives their degrees and largest eigenvalue
multiplicities on every noncentral class.  The columns are the
exact class indices of the retained CTblLib input; the order row
prevents a class-name convention from being tacitly substituted.
\[
\begin{array}{c|r|rrrrrrrrrrr}
\text{row}&\deg&3&4&5&6&7&8&9&10&11&12&13\\\hline
\text{order}&&4&3&6&3&6&8&8&5&10&5&10\\\hline
8 &4 &2&2&2&2&2&1&1&1&1&1&1\\
9 &4 &2&2&2&2&2&1&1&1&1&1&1\\
10&8 &4&3&3&3&3&2&2&2&2&2&2\\
11&8 &4&3&3&3&3&2&2&2&2&2&2\\
12&10&5&4&4&4&4&3&3&2&2&2&2\\
13&10&5&4&4&4&4&3&3&2&2&2&2
\end{array}                                                   \tag{SPINQ.2}
\]
Every entry is at most half its row degree.  For completeness,
the exact calculation is the cyclic Fourier formula: if $h$ has
order $m$ and $\zeta_m$ is a primitive $m$th root, then the
multiplicity of $\zeta_m^k$ is
\[
       \frac1m\sum_{r=0}^{m-1}\chi(h^r)\zeta_m^{-kr}.           \tag{SPINQ.3}
\]
The retained complete character table and power maps supply all
terms.  Prime power maps determine $h^d$ for $d=\gcd(r,m)$;
the remaining coprime power is the corresponding cyclotomic
Galois automorphism on its character value.  Reducing the sum
in $\mathbf Z[x]/(\Phi_m)$ gives a constant divisible by $m$.
The exact script checks nonnegativity, integrality and sum equal
to the degree for each of the $66$ spectra, then gives (SPINQ.2).
It selects the six rows from their actual central-involution
values, rather than assuming a supplied row-selection list.
The full published table is an explicit cited input
\cite{SPOR:CTblLib}; this calculation does not purport to reconstruct
the table or identify an unknown group from numerical orthogonality.
\end{proof}

\begin{theorem}[HD-SPIN-02b: all faithful spin normal cores]
\label{HD-SPIN-02b}
Let a finite-cover sixfold source have an inducing image $J$
with an irreducibly acting normal central product
\[
                  E=E_1\cdots E_t,\qquad
                  E_i\simeq2.\mathfrak A_{n_i},\quad n_i\ge8,
                                                               \tag{SPINQ.4}
\]
where the factor representations are faithful.  Then the actual
six-dimensional multiplicity variation has period-map rank at
most eight.  This holds for either polarization determinant class,
any inducing character field, and all compatible outer
automorphisms and permutations of tensor factors.
\end{theorem}

\begin{proof}
First take a single factor $N=2.\mathfrak A_n$ on a faithful
irreducible space $W$ of dimension $a$, and a nonscalar element
$g$ of its linear normalizer with finite order up to scalar.
The complete group-theoretic argument of HD-SPIN-01a gives a
lift $t$ of a three-cycle for which $h=[g,t]$ is noncentral in
a natural $2.\mathfrak A_6$ subgroup.  The central involution
acts as $-1$ on every restriction constituent.  HD-SPIN-02a
therefore gives $\nu(h)\ge a/2$.  The same exact commutator
rank inequality, with $g$ rescaled at a largest eigenvalue, gives
\[
                  a/2\le\nu(h)\le2\nu(g),
                  \qquad \nu(g)\ge a/4.                    \tag{SPINQ.5}
\]
The restriction can be reducible: (SPINQ.1) holds for each fixed
eigenvalue across all constituents, which is precisely what is
needed here.

Write $V=\bigotimes_iV_i$ for (SPINQ.4).  Every $a_i=\dim V_i$
is at least eight by the published general spin minimum used in
(SPIN.5).  The intrinsic factor-algebra and tensor-normalizer
argument in D-SPOR-03 applies.  A factor-preserving nonscalar
$g$ satisfies $\nu(g)\ge\dim V/4$ by (SPINQ.5) and the tensor
eigenspace inequality.  A nontrivial factor permutation satisfies
the same bound by the weighted cyclic-word count (SPOR.6).
Consequently every nonscalar element satisfies this quarter bound,
and every genuine scalar inertia moves the entire space.

Put $e=\dim V\ge8$.  The intermediate-cover and period-factorization
argument in D-SPOR-03 applies without any sporadic hypothesis.
Explicitly, if its base has genus $h$, the actual six-dimensional
multiplicity cohomology gives
\[
          6=2e(h-1)+\sum_j\operatorname{rank}(1-g_j).
\]
For $h\ge2$ this is impossible.  For $h=1$ there are at most six
genuine branch points and its pointed-curve moduli, hence period
image, has dimension at most six.  For $h=0$, period rank nine
would require at least twelve genuine branch points.  The
quarter bound would then imply
\[
          3e\le12\lceil e/4\rceil
             \le\sum_j\operatorname{rank}(1-g_j)=2e+6,
\]
so $e\le6$, contrary to $e\ge8$.  Therefore rank nine is
impossible, as asserted.
\end{proof}

\paragraph{Checks and precise remaining scope.}
The complete local spectra and source-input hash are retained in
\texttt{research/spin\_quarter\_certificate.json}; its integer-polynomial
implementation is separate from the peer's Ramanujan-trace
implementation.  The theorem supersedes the ESD restriction in
HD-SPIN-01b.  Ordinary alternating representations and exceptional
central covers in alternating degrees below eight remain outside
this result.  No new detector cycle, coefficient-field descent,
or actual dominant component is inferred from this exclusion.
